% Test verisi: PDF'te aramanin satir sonu kapisinin kaynagi.
%
% NEDEN VAR. pdfium satir sonunda tireyle bolunmus sozcugu
% `karsilastirilabi\x02lirlik` diye (U+0002, satir sonu koymadan), satir
% sonunu da bosluk yerine `\r\n` diye veriyor. Arama ikisini de bulamiyordu
% (olcum `gui.pdf_metin.arama_metni`de). Kapi GERCEK pdfium ciktisina
% bakiyor: pdfium bu bicimi degistirirse kapi da duser.
%
% Dar sutun ve uzun sozcukler bolunmeyi garantiliyor; TEK sayfa, cunku
% arama sayfa basina yapiliyor ve sayfa sonunda bolunen sozcuk bu kapinin
% konusu degil. Govdede yalniz duz sozcuk ve nokta var: kapi beklenen
% gecisleri bu dosyanin govdesinden sayiyor.
%
% YENIDEN URETMEK ICIN (T1 + Turkce babel, Turkce tezlerin olagan duzeni):
%     pdflatex -output-directory=tests/veri tests/veri/tireli_arama.tex
% (yardimci dosyalar .gitignore'da; depoda .tex ve .pdf duruyor.)
\documentclass{article}
\usepackage[T1]{fontenc}
\usepackage[turkish]{babel}
\usepackage[paperwidth=10cm,paperheight=24cm,textwidth=5cm,top=1cm,bottom=1cm]{geometry}
\pagestyle{empty}
\begin{document}
Karşılaştırılabilirlik ve sürdürülebilirlik bu tezin iki temel ölçütüdür.
Uygulanabilirliğinin değerlendirilmesinde istatistiksel yöntemler kullanıldı.

Yapılandırılmış veriler bilgisayarlaştırılmış bir sistemde toplandı ve
parametrelendirilmiş modeller gerçekleştirilmesi zor deneylerde sınandı.

Çözümlenebilirlik hesaplanabilirlik ile birlikte ele alındı. Coğrafi
fotoğraflar ve fiziksel ölçümler üniversite laboratuvarında incelendi.

Öğrencilerin katılımıyla yürütülen çalışmada karşılaştırılabilirlik yeniden
değerlendirildi ve sürdürülebilirlik için yeni bir yaklaşım geliştirildi.
\end{document}
